% ===== PRÉAMBULE CANONIQUE — à coller VERBATIM en tête de CHAQUE .tex =====
\documentclass[11pt,a4paper]{article}
\usepackage[utf8]{inputenc}\usepackage[T1]{fontenc}\usepackage{lmodern}
\usepackage[french]{babel}
\usepackage{amsmath,amssymb,mathtools}\usepackage{icomma}
\usepackage[version=4]{mhchem}          % \ce{...} : équations & formules
\usepackage{siunitx}\sisetup{locale=FR} % \qty{}{}, \si{}, \ang{}
\usepackage{chemfig}                    % \chemfig{...} : structures développées
\usepackage{xcolor}\usepackage{colortbl}
\usepackage{tikz}\usepackage{pgfplots}\pgfplotsset{compat=1.18}
\usepackage[most]{tcolorbox}\usepackage{enumitem}\usepackage{array,booktabs}
\usepackage[a4paper,margin=2.2cm]{geometry}
\usetikzlibrary{arrows.meta,calc,angles,quotes,positioning,patterns,backgrounds,shapes.geometric}
\definecolor{primary}{RGB}{222,49,99}\definecolor{primarydark}{RGB}{150,22,55}
\definecolor{defcol}{RGB}{0,114,178}\definecolor{methcol}{RGB}{52,92,148}
\definecolor{excol}{RGB}{0,135,90}\definecolor{attncol}{RGB}{200,40,55}
\tcbset{boxbase/.style={enhanced,breakable,boxrule=0.9pt,arc=2.5pt,
  left=3mm,right=3mm,top=2.3mm,bottom=2.3mm,fonttitle=\bfseries,coltitle=white}}
% --- boîtes TUTORIEL (noms EXACTS, ne pas renommer) ---
\newtcolorbox{cle}[1][]{boxbase,colback=defcol!6,colframe=defcol,
  title={\textsc{À retenir}\ifx\relax#1\relax\relax\else\textmd{ : \itshape #1}\fi}}
\newtcolorbox{methode}[1][]{boxbase,colback=methcol!7,colframe=methcol,sharp corners,
  title={\textsc{Méthode}\ifx\relax#1\relax\relax\else\textmd{ --- \itshape #1}\fi}}
\newtcolorbox{exemplebox}[1][]{boxbase,colback=excol!5,colframe=excol,
  title={\textsc{Exemple}\ifx\relax#1\relax\relax\else\textmd{ : \itshape #1}\fi}}
\newtcolorbox{piege}[1][]{boxbase,colback=attncol!6,colframe=attncol,
  title={\textsc{Piège}\ifx\relax#1\relax\relax\else\textmd{ : \itshape #1}\fi}}
\newtcolorbox{remarque}[1][]{boxbase,colback=black!4,colframe=black!45,
  title={\textsc{Remarque}\ifx\relax#1\relax\relax\else\textmd{ : \itshape #1}\fi}}
% --- boîtes EXERCICE / CORRIGÉ (noms EXACTS, auto-comptées) ---
\newtcolorbox[auto counter]{exo}[1][]{boxbase,colback=primary!4,colframe=primary,
  title={\textsc{Exercice}~\thetcbcounter\ifx\relax#1\relax\relax\else\textmd{ --- \itshape #1}\fi}}
\newtcolorbox[auto counter]{corrige}[1][]{boxbase,colback=excol!4,colframe=excol,
  title={\textsc{Corrigé}~\thetcbcounter\ifx\relax#1\relax\relax\else\textmd{ --- \itshape #1}\fi}}
\newcommand{\R}{\mathbb{R}}\newcommand{\N}{\mathbb{N}}\newcommand{\Z}{\mathbb{Z}}\newcommand{\ds}{\displaystyle}
% (un \usepackage{fancyhdr}+en-tête est possible ; il est ignoré côté web)
% =========================================================================
\begin{document}
\begin{exo}[combustions et molécules oxygénées]
Pondère chaque combustion complète (produits : \ce{CO2} et \ce{H2O}). Attention à l'oxygène déjà
présent dans certains réactifs.
\begin{enumerate}[label=\alph*)]
  \item \ce{C4H10 + O2 -> CO2 + H2O}
  \item \ce{C2H5OH + O2 -> CO2 + H2O}
  \item \ce{C6H12O6 + O2 -> CO2 + H2O}
\end{enumerate}
\end{exo}

\begin{exo}[groupements polyatomiques conservés]
Pondère en traitant chaque groupement (\ce{OH}, \ce{SO4}\ldots) comme un bloc.
\begin{enumerate}[label=\alph*)]
  \item \ce{Al(OH)3 + H2SO4 -> Al2(SO4)3 + H2O}
  \item \ce{Ca(OH)2 + HCl -> CaCl2 + H2O}
  \item \ce{Fe2O3 + HCl -> FeCl3 + H2O}
\end{enumerate}
\end{exo}

\begin{exo}[réactions de précipitation]
Pondère chaque réaction de double échange.
\begin{enumerate}[label=\alph*)]
  \item \ce{AgNO3 + CaCl2 -> AgCl v + Ca(NO3)2}
  \item \ce{Pb(NO3)2 + KI -> PbI2 v + KNO3}
  \item \ce{Na3PO4 + CaCl2 -> Ca3(PO4)2 v + NaCl}
\end{enumerate}
\end{exo}

\begin{exo}[azote, soufre et gros coefficients]
Pondère chaque équation.
\begin{enumerate}[label=\alph*)]
  \item \ce{NH3 + O2 -> NO + H2O}
  \item \ce{C8H18 + O2 -> CO2 + H2O}
  \item \ce{FeS2 + O2 -> Fe2O3 + SO2}
\end{enumerate}
\end{exo}

\begin{exo}[vérifier, corriger, sommer]
Un étudiant propose l'équation « ajustée » suivante pour l'attaque de l'aluminium par l'acide
chlorhydrique :
\[
\ce{Al + 3 HCl -> AlCl3 + H2 ^}.
\]
\begin{enumerate}[label=\alph*)]
  \item Est-elle correctement pondérée ? Justifie par l'inventaire des atomes.
  \item Si non, écris l'équation correctement pondérée (plus petits entiers).
  \item Donne la somme des coefficients stœchiométriques de l'équation correcte.
\end{enumerate}
\end{exo}

\end{document}
